European option contracts specify an expiration date, and if the option is to be exercised at all, the exercise must occur on the expiration date. An option whose owner can choose to exercise at any time up to and including the expiration date is called \textit{American}.

Whereas a European option can be exercised only at a fixed date, an American option can be exercised any time up to its expiration. The value of an American option is the value achieved by exercising optimally. Finding this value entails finding the optimal exercise rule by solving an optimal stopping problem and computing the expected discounted payoff of the option under this rule. The embedded optimisation problem makes this a difficult problem for simulation.

Because an American option can be exercised at any time prior to its expiration,
it can never be worth less than the payoff associated with immediate exercise.
This is called the \textit{intrinsic value} of the option.

To price an American option, just as with a European option, we could imagine
selling the option in exchange for some initial capital and then consider how to
use this capital to hedge the short position in the option. In this case, we
would need to be ready to pay off the option at all times prior to the
expiration date because we do not know when it will be exercised. We could
determine when, from our point of view, is the worst time for the owner to
exercise the option. From the owner's point of view, this would be the \textit{optimal
exercise time}, and we shall call it that. We could then compute the initial
capital we need in order to be hedged against exercise at the optimal exercise
time. Finally, we could show how to invest this capital so that we are hedged
even if the owner exercises at a nonoptimal time. In the subsequent sections, we
do all these things but begin the analysis at a different point than for
European options. We define the price of American options using a risk-neutral
pricing formula and then show that this price is the smallest initial capital
that permits construction of the hedge just described.

\subsection{Perpetual American Put}

The simplest interesting American option is the \textit{perpetual American put}. It is interesting because the optimal exercise policy is not obvious, and it is simple because this policy can be determined explicitly. Although this is not a traded option, we begin our discussion with it in order to present in a simple context the ideas behind the subsequent analysis of more realistic options.
The underlying asset in most of this chapter (except where the asset pays dividends, which will be explicitly mentioned) has the price process $S(t)$ given by
\begin{equation}
	\label{eq:GBMRiskNeutral}
  dS(t) = rS(t) , dt + \sigma S(t) , d\widetilde{W}(t), \tag{8.3.1}
\end{equation}
where the interest rate $r$ and the volatility $\sigma$ are strictly positive constants and $\widetilde{W}(t)$ is a Brownian motion under the risk-neutral probability measure $\widetilde{\Prob}$. The perpetual American put pays $K - S(t)$ if it is exercised at time $t$. This is its intrinsic value.

\begin{defn}{Perpetual American Put}{PerpetualAmericanPut}
  Let $\mathcal{T}$ be the set of all stopping times. The price of the perpetual American put is defined to be
  \begin{equation}
	\label{eq:PerpetualAmericanPutPrice}
    v_*(x) = \max_{\tau \in \mathcal{T}} \widetilde{\E} \left[ e^{-r\tau} (K - S(\tau)) \right],
  \end{equation}
  where $x = S(0)$ in (8.3.2) is the initial stock price. In the event that $\tau = \infty$, we interpret $e^{-r\tau}(K - S(\tau))$ to be zero.
\end{defn}

The idea behind \Cref{defn:PerpetualAmericanPut} is that the owner of the perpetual American put can choose an exercise time $\tau$, subject only to the condition that she may not look ahead to determine when to exercise. The mathematical formulation of this ``not look ahead'' restriction is that $\tau$ must be a stopping time. The price of the option at time zero is the risk-neutral expected payoff of the option, discounted from the exercise time back to time zero. If the option is never exercised, its payoff is zero. This explains the term under the expectation on the right-hand side of \Cref{eq:PerpetualAmericanPutPrice}. The owner of the option should choose the exercise strategy that maximizes this expected payoff, discounted back to time zero, and thus we define the price of the option to be the maximum over $\tau \in \mathcal{T}$ of the discounted expected payoffs.

The owner of the perpetual American put can exercise at any time. In particular, there is no expiration date after which the put can no longer be exercised. This makes every date like every other date; the time remaining to expiration is always the same (i.e., infinity). Because every date is like every other date, it is reasonable to expect that the optimal exercise policy depends only on the value of $S(t)$ and not on the time variable $t$. The owner of the put should exercise as soon as $S(t)$ falls far enough below $K$. In other words, it is reasonable to expect that the optimal exercise policy is of the form
\begin{quote} Exercise the put as soon as $S(t)$ falls to the level $L_*$.
\end{quote}

We have two questions to answer:
\begin{enumerate}[label=(\roman*)] \item What is the value of $L_*$, and how do we know it corresponds to optimal exercise?
\item What is the value of the put?
\end{enumerate}

For the perpetual American put, we can base the answers to these questions on explicit computations.


\subsection{Approximate Continuation Values}
Regression-based methods posit an expression for the continuation value of the form
\begin{equation} \label{eq:ContinuationValueExpectation}
	\E[V_{i+1}(X_{i+1}) \mid X_i = x] = \sum_{r=1}^M \beta_{ir} \psi_r(x),
\end{equation}
for some basis functions $\psi_r : \R^d \to \R$ and constants $\beta_{ir}, r = 1, \dots, M$. Using the notation $C_i$ for the continuation value at time $i$, we may equivalently write \Cref{eq:ContinuationValueExpectation} as
\begin{equation} \label{eq:ContinuationValueVectorForm}
	C_i(x) = \beta_i^\top \psi(x),
\end{equation}
with
\begin{equation*}
	\beta_i^\top = (\beta_{i1}, \dots, \beta_{iM}), \quad \psi(x) = (\psi_1(x), \dots, \psi_M(x))^\top.
\end{equation*}
One could use different basis functions at different exercise dates, but to simplify notation, we suppress any dependence of $\psi$ on $i$.

Assuming a relation of the form \Cref{eq:ContinuationValueExpectation} holds, the vector $\beta_i$ is given by
\begin{equation} \label{eq:BetaVectorDefinition}
	\beta_i = (\E[\psi(X_i)\psi(X_i)^\top])^{-1} \E[\psi(X_i)V_{i+1}(X_{i+1})] \equiv B_\psi^{-1} B_{\psi V}.
\end{equation}
Here, $B_\psi$ is the indicated $M \times M$ matrix (assumed nonsingular) and $B_{\psi V}$ the indicated vector of length $M$. The variables $(X_i, X_{i+1})$ inside the expectations have a joint distribution corresponding to the state of the underlying Markov chain at dates $i$ and $i+1$.

The coefficients $\beta_{ir}$ could be estimated from observations of pairs $(X_{ij}, V_{i+1}(X_{i+1,j}))$, $j = 1, \dots, b$, each consisting of the state at time $i$ and the corresponding option value at time $i+1$. Consider, in particular, independent paths $(X_{1j}, \dots, X_{mj})$, $j = 1, \dots, b$, and suppose for a moment that the values $V_{i+1}(X_{i+1,j})$ are known. The least-squares estimate of $\beta_i$ is then given by
\begin{equation*}
	\hat{\beta}_i = \hat{B}_\psi^{-1} \hat{B}_{\psi V},
\end{equation*}
where $\hat{B}_\psi$ and $\hat{B}_{\psi V}$ are the sample counterparts of $B_\psi$ and $B_{\psi V}$. More explicitly, $\hat{B}_\psi$ is the $M \times M$ matrix with $qr$ entry
\begin{equation*}
	\frac{1}{b} \sum_{j=1}^b \psi_q(X_{ij}) \psi_r(X_{ij})
\end{equation*}
and $\hat{B}_{\psi V}$ is the $M$-vector with $r$th entry
\begin{equation} \label{eq:SampleCounterpartVector}
	\frac{1}{b} \sum_{k=1}^b \psi_r(X_{ik}) V_{i+1}(X_{i+1, k}).
\end{equation}
All of these quantities can be calculated from function values at pairs of consecutive nodes $(X_{ij}, X_{i+1, j})$, $j = 1, \dots, b$. In practice, $V_{i+1}$ is unknown and must be replaced by estimated values $\hat{V}_{i+1}$ at downstream nodes. The estimate $\hat{\beta}_i$ then defines an estimate
\begin{equation} \label{eq:EstimatedContinuationValue}
	\hat{C}_i(x) = \hat{\beta}_i^\top \psi(x),
\end{equation}
of the continuation value at an arbitrary point $x$ in the state space $\R^d$. This, in turn, defines a procedure for estimating the option value.

\begin{algorithm}[H]
	\caption{Regression-Based Pricing Algorithm}
	\label{alg:GeneralRegressionPricing}
	\begin{algorithmic}[1]
		\State Simulate $b$ independent paths $\{X_{1j}, \dots, X_{mj}\}, j = 1, \dots, b$, of the Markov chain.
		\State At terminal nodes, set $\hat{V}_{mj} = h_m(X_{mj})$ for $j = 1, \dots, b$.
		\For{$i = m-1$ \textbf{down to} $1$}
		\State Given estimated values $\hat{V}_{i+1, j}, j = 1, \dots, b$, use regression to calculate $\hat{\beta}_i = \hat{B}_\psi^{-1} \hat{B}_{\psi V}$.
		\State Compute estimated continuation values $\hat{C}_i(X_{ij})$ using \Cref{eq:EstimatedContinuationValue}.
		\For{$j = 1, \dots, b$}
		\State Set value based on immediate exercise versus estimated continuation:
		\begin{equation} \label{eq:GeneralRegressionUpdate}
			\hat{V}_{ij} = \max \left\{ h_i(X_{ij}), \hat{C}_i(X_{ij}) \right\}.
		\end{equation}
		\EndFor
		\EndFor
		\State Set $\hat{V}_0 = (\hat{V}_{11} + \dots + \hat{V}_{1b}) / b$.
	\end{algorithmic}
\end{algorithm}

\begin{defn}{American Contingent Claim}{AmericanContingentClaim}
	An \textit{American contingent claim} $X^a = (X, \mathcal{T}_{[0,T]})$ expiring at time $T$ consists of a sequence of payoffs $(X_t)_{t=0}^T$, where $X_t$ is an $\cF_t$-measurable random variable for $t = 0, 1, \dots, T$, and of the set $\mathcal{T}_{[0,T]}$ of admissible exercise policies.
\end{defn}

\begin{prop}{Arbitrage Price of American Put}{ArbitragePriceAmericanPut}
	The arbitrage price $P_t^a$ of an American put option equals
	\begin{equation} \label{eq:AmericanPutArbitragePrice}
		P_t^a = \max_{\tau \in \mathcal{T}_{[t,T]}} \E_{\Pm^*} \left( \hat{r}^{-(\tau-t)} (K - S_\tau)^+ \mid \cF_t \right), \quad \forall t \le T.
	\end{equation}
	Moreover, for any $t \le T$ the stopping time $\tau_t^*$ that realizes the maximum in \Cref{eq:AmericanPutArbitragePrice} is given by the expression (by convention $\min \emptyset = T$)
	\begin{equation} \label{eq:OptimalStoppingTime}
		\tau_t^* = \min \left\{ u \in \{t, t + 1, \dots, T\} \mid (K - S_u)^+ \ge P_u^a \right\}.
	\end{equation}
\end{prop}

\begin{prop}{Arbitrage Price of American Claim in CRR}{ArbitragePriceAmericanClaimCRR}
	For every $t \le T$, the arbitrage price $\pi(X^a)$ of an arbitrary American claim $X^a$ in the CRR model equals
	\begin{equation} \label{eq:GeneralAmericanArbitragePrice}
		\pi_t(X^a) = \max_{\tau \in \mathcal{T}_{[t,T]}} \E_{\Pm^*} \left( \hat{r}^{-(\tau-t)} X_\tau \mid \cF_t \right).
	\end{equation}
	The price process $\pi(X^a)$ can be determined using the following recurrence relation, for every $t \le T - 1$,
	\begin{equation} \label{eq:SnellEnvelopeRecurrence}
		\pi_t(X^a) = \max \left\{ X_t, \E_{\Pm^*} \left( \hat{r}^{-1} \pi_{t+1}(X^a) \mid \cF_t \right) \right\}
	\end{equation}
	subject to the terminal condition $\pi_T(X^a) = X_T$. In the case of a path-independent American claim $X^a$ with the reward function $h$ we have that, for every $t \le T - 1$,
	\begin{equation} \label{eq:PathIndependentRecurrence}
		\pi_t(X^a) = \max \left\{ h(S_t, t), \hat{r}^{-1} \left( p_* \pi_{t+1}^u(X^a) + (1 - p_*) \pi_{t+1}^d(X^a) \right) \right\},
	\end{equation}
	where, for a generic stock price value $S_t$ at time $t$, we write $\pi_{t+1}^u(X^a)$ and $\pi_{t+1}^d(X^a)$ to denote the values of the price process $\pi(X^a)$ at time $t + 1$ in the nodes that correspond to the upward and downward movements of the stock price during the time-period $[t, t + 1]$, that is, for the values $uS_t$ and $dS_t$ of the stock price at time $t + 1$, respectively.
\end{prop}
